\subsection{有理数域}

定义4. 整数环Z的分式域称为有理数域，记作Q。Q的元素称为有理数。

因此，每个有理数都具有形式 \(a / b\) ，其中 \(a\) 与 \(b\) 为满足 \(b \neq 0\) 的有理整数（我们甚至可以取 \(b > 0\) ，这由关系

\[
a / b = (- a) / (- b)
\]

所证明）。 \(\mathbf{Q}_{+}\) 用来表示形如 \(a / b\) （其中 \(a\in \mathbf{N}\) ， \(b\in \mathbf{N}^*\) ）的有理数所成的集合。

我们有如下关系：

\begin{enumerate}
\def\labelenumi{(\arabic{enumi})}
\tightlist
\item
\end{enumerate}

\[
\mathbf {Q} _ {+} + \mathbf {Q} _ {+} = \mathbf {Q} _ {+}\tag{2}
\]

\[
\mathbf {Q} _ {+}. \mathbf {Q} _ {+} = \mathbf {Q} _ {+}\tag{3}
\]

\[
\mathbf {Q} _ {+} \cap (- \mathbf {Q} _ {+}) = \{0 \}\tag{4}
\]

\[
\mathbf {Q} _ {+} \cup (- \mathbf {Q} _ {+}) = \mathbf {Q}\tag{5}
\]

\[
\mathbf {Q} _ {+} \cap \mathbf {Z} = \mathbf {N}.
\]

前两个关系由公式 \(a / b + a' / b' = (ab' + a'b) / bb'\) , \((a / b)(a'/b') = aa'/bb'\) , \(0 \in \mathbf{Q}_+\) , \(1 \in \mathbf{Q}_+\) 以及 \(\mathbf{N}\) 对加法和乘法封闭、\(\mathbf{N}^*\) 对乘法封闭这一事实得出。显然 \(0 \in \mathbf{Q}_+\) ，从而 \(0 \in (-\mathbf{Q}_+)\) ；设 \(x\) 属于 \(\mathbf{Q}_+ \cap (-\mathbf{Q}_+)\) ；则存在正整数 \(a, b, a', b'\) 满足 \(b \neq 0\) , \(b' \neq 0\) 且 \(x = a / b = -a'/b'\) ；于是 \(ab' + ba' = 0\) ，从而 \(ab' = 0\) （因为 \(ab' \geqslant 0\) 且 \(ba' \geqslant 0\) ），进而得 \(a = 0\) ，因为 \(b' \neq 0\) ；换言之， \(x = 0\) 。最后，显然 \(\mathbf{N} \subset \mathbf{Z}\) 且 \(\mathbf{N} \subset \mathbf{Q}_+\) 。反之，若 \(x\) 属于 \(\mathbf{Z} \cap \mathbf{Q}_+\) ，则它是一个有理整数；存在两个有理整数 \(a\) 和 \(b\) 满足 \(a \geqslant 0\) , \(b > 0\) 且 \(x = a / b\) ，于是 \(a = bx\) ；若 \(x \notin \mathbf{N}\) ，则 \(-x > 0\) ，于是 \(-a = b(-x) > 0\) ，从而 \(a < 0\) ，与假设矛盾。

给定两个有理数 \(x\) 和 \(y\) ，我们记 \(x \leqslant y\) ，若 \(y - x \in \mathbf{Q}_+\) 。由公式(1)、(3)和(4)容易推出 \(x \leqslant y\) 是

Q上的一个全序，由(5)可知该关系在Z上诱导出通常的序关系。最后，由(1)可得关系 \(x \leqslant y\) 和 \(x' \leqslant y'\) 蕴含 \(x + x' \leqslant y + y'\) ，由(2)可得关系 \(x \leqslant y\) 蕴含 \(xz \leqslant yz\) 对一切 \(z \geqslant 0\) ，且蕴含 \(xz \geqslant yz\) 对 \(z \leqslant 0\) （参见VI，§2，第1号）。

设 \(x\) 是一个有理数。称 \(x\) 为正的，若 \(x \geqslant 0\) ；称其为严格正的，若 \(x > 0\) ；称其为负的，若 \(x \leqslant 0\) ；称其为严格负的，若 \(x < 0\) 。正有理数的集合是 \(\mathbf{Q}_{+}\) ，负有理数的集合是 \(-\mathbf{Q}_{+}\) 。若 \(\mathbf{Q}^{*}\) 表示非零有理数的集合，则严格正数的集合 \(\mathbf{Q}_{+}^{*}\) 等于 \(\mathbf{Q}^{*} \cap \mathbf{Q}_{+}\) ，而 \(-\mathbf{Q}_{+}^{*}\) 是严格负数的集合。

集合 \(\mathbf{Q}_{+}^{*}\) 和 \(\{1, -1\}\) 是乘法群 \(\mathbf{Q}^{*}\) 的子群。每个有理数 \(x \neq 0\) 都能以且仅以一种方式表示成 \(1.y\) 或 \((-1).y\) 的形式，其中 \(y > 0\) ；因此乘法群 \(\mathbf{Q}^{*}\) 是子群 \(\mathbf{Q}_{+}^{*}\) 与 \(\{1, -1\}\) 的乘积， \(x\) 在 \(\mathbf{Q}_{+}^{*}\) 中的分量称为 \(x\) 的绝对值，记作 \(|x|\) ； \(x\) 在 \(\{-1, 1\}\) 中的分量（等于 1 若 \(x > 0\) ，等于 \(-1\) 若 \(x < 0\) ）称为 \(x\) 的符号，记作 \(\operatorname{sgn} x\) .

通常，我们把这两个函数延拓到整个 \(\mathbf{Q}\) 上，办法是设定 \(|0| = 0\) 和 \(\operatorname{sgn} 0 = 0\) 。

